
\begin{proposition}{7.4.56}\label{prop:7-4-56}
有
\begin{equation}
\begin{aligned}
&\sum_i(-1)^{*(i)}
\mathfrak g\bigl(
Y_1,\ldots,\delta(Y_i),\ldots,Y_\ell
\bigr)(\mathbf x)\\
&\quad+
\sum_{i<j}(-1)^{*(i,j)}
\mathfrak g\bigl(
[Y_i,Y_j],Y_1,\ldots,
\widehat Y_i,\ldots,\widehat Y_j,
\ldots,Y_\ell\bigr)(\mathbf x)\\
&=
\sum_c
\sum_{\mathbf Y_{(1)}
\cup\mathbf Y_{(2)}=\mathbf Y}
(-1)^{*(\mathbf Y_{(1)},\mathbf Y_{(2)},c)}
\\[-0.3ex]
&\qquad
\mathfrak g\bigl(\mathbf Y_{(1)}\bigr)
\bigl(\mathbf x_c^{3;1},
\mathfrak g(\mathbf Y_{(2)})
(\mathbf x_c^{3;2}),
\mathbf x_c^{3;3}\bigr).
\end{aligned}
\tag{7.4.57}\label{eq:7-4-57}
\end{equation}
\end{proposition}

右端第一层求和中的
$\mathbf Y_{(1)},\mathbf Y_{(2)}$
取遍 $\mathbf Y$ 的所有洗牌：
二者均为 $\mathbf Y$ 的保序子列，
彼此不交，且其并为 $\mathbf Y$。
式~\eqref{eq:7-4-57}
中的符号指数规定如下：
\[
\begin{aligned}
*(i)
&=\deg'Y_1+\cdots+\deg'Y_{i-1}+1,\\
*(i,j)
&=\deg'Y_i\sum_{k<i}\deg'Y_k
+\deg'Y_j\sum_{\substack{k<j\\k\ne i}}
\deg'Y_k+\deg'Y_i,\\
*(\mathbf Y_{(1)},\mathbf Y_{(2)},c)
&=\deg'\mathbf x_c^{(3;1)}
\deg'\mathbf Y_{(2)}
+\deg'\mathbf x_c^{(3;1)}
+\deg'\mathbf Y_{(1)}\\
&\quad+
\sum_{\substack{a\in\mathbf Y_{(1)},\,
b\in\mathbf Y_{(2)}\\b<a}}
\deg'Y_a\deg'Y_b+1.
\end{aligned}
\]
注意
$\deg'\mathbf Y
=\sum_i(\deg Y_i-1)$。

为证明命题~\ref{prop:7-4-56}，
我们先引入更多记号，
并将与树关联的同态构造
推广到带附加结构装饰的情形。
考虑一对
$\Gamma=(|\Gamma|,Z)$，
其中 $|\Gamma|$ 为平面树，
$Z$ 是附加结构。
设 $(D,\delta)$ 为
$\mathbb Q$ 上任意上链复形。
\begin{enumerate}
\item[\textup{(7.4.58.1)}]\label{cond:7-4-58-1}
将 $|\Gamma|$ 的顶点集
$C^0(\Gamma)$ 分成不交并
\[
C^0(\Gamma)
=C^0_{\mathrm{int}}(\Gamma)
\cup C^0_{\mathrm{ext}}(\Gamma).
\]
这两部分分别为内顶点集与外顶点集，
并满足下列条件。

\item[\textup{(7.4.58.2)}]\label{cond:7-4-58-2}
每个外顶点恰邻接一条边，
称为外边；其余边称为内边。

\item[\textup{(7.4.58.3)}]\label{cond:7-4-58-3}
外顶点恰有 $k+1$ 个，
编号为 $0,1,\ldots,k$。
这一编号遵循平面逆时针定向
诱导的循环次序。\liftIndex{I291}

\item[\textup{(7.4.58.4)}]\label{cond:7-4-58-4}
若内顶点 $v$ 邻接 $k_v+1$ 条边，
则赋予次数为 $d_v$ 的同态
$Z(v)\in\operatorname{Hom}(B_{k_v}D,D)$。

\item[\textup{(7.4.58.5)}]\label{cond:7-4-58-5}
若 $e$ 是一条边，
则赋予次数为 $d_e$ 的同态
$Z(e)\in\operatorname{Hom}(D,D)$。
\end{enumerate}

对每个
$\Gamma=(|\Gamma|,Z)$，
可按树的关联关系复合
$Z(v)$ 与 $Z(e)$，
得到次数为
$\sum_vd_v+\sum_ed_e$
的同态
\[
\mathfrak h(\Gamma)
\in\operatorname{Hom}(B_kD,D).
\]
然而，复合中的符号约定
颇为微妙；为确定它，
我们按内顶点数
归纳定义算子 $\mathfrak h(\Gamma)$。

先设 $|\Gamma|$
只有一个内顶点 $v$，
其邻接的边为
$e_0,\ldots,e_k$。
对 $i=0,\ldots,k$，
$e_i$ 的另一端点
是第 $i$ 个外顶点。
定义
\begin{equation}
\mathfrak h(\Gamma)
=Z(e_0)\circ Z(v)\circ
\bigl(
Z(e_1)\widehat\otimes\cdots
\widehat\otimes Z(e_k)
\bigr).
\tag{7.4.59}\label{eq:7-4-59}
\end{equation}
这里 $\widehat\otimes$ 定义为
\[
\begin{aligned}
&(A_1\widehat\otimes\cdots
\widehat\otimes A_n)
(u_1\otimes\cdots\otimes u_n)\\
&\qquad=(-1)^*
A_1(u_1)\otimes\cdots\otimes A_n(u_n),
\qquad
*=\sum_{i<j}\deg A_i\,\deg u_j.
\end{aligned}
\]
它与式~\eqref{eq:7-4-53}
定义的 $\widehat\otimes'$
略有不同：这里不用移位次数。

再设 $\Gamma$ 至少有两个内顶点\translatornote{源书此处只写“至少两个顶点”；但单内顶点树已有外顶点，故按内顶点数归纳须这样限定。见疑误 A-0570。}。
取含第 $0$ 个外顶点 $v_0$ 的边
$e_0$（即第 $0$ 条外边），
记
$\partial e_0=\{v_0,v'_0\}$。
从 $\Gamma$ 中移去 $e_0$
及其两个端点，
如前得到
$\Gamma_1,\ldots,\Gamma_n$。
按定义，
每棵 $\Gamma_i$
（$i=1,\ldots,n$）
的内顶点数均少于 $\Gamma$\translatornote{源书此处将下标范围写成 $i=1,\ldots,k$，但前后均为 $n$ 棵子树，故译文据此改用 $n$。见疑误 A-0569。}。
因此可用公式
\[
\mathfrak h(\Gamma)
=Z(e_0)\circ Z(v'_0)\circ
\bigl(
\mathfrak h(\Gamma_1)\widehat\otimes
\cdots\widehat\otimes
\mathfrak h(\Gamma_n)
\bigr)
\]
归纳定义 $\mathfrak h(\Gamma)$。

上述符号规则也可这样描述。
作为例子，
考虑图~\ref{fig:7-4-5}
中的图 $\Gamma$。
\par\medskip\noindent\begin{minipage}{\textwidth}
\centering
\setlength{\unitlength}{0.65mm}
\begin{picture}(105,69)
\linethickness{0.55pt}
\put(56,58){\line(0,-1){13}}
\qbezier(56,45)(36.5,24)(17,3)
\qbezier(56,45)(65.5,32.5)(75,20)
\qbezier(75,20)(68.5,11.5)(62,3)
\qbezier(75,20)(81.5,11.5)(88,3)
\put(53,61){$\scriptstyle v_0$}
\put(59,51){$\scriptstyle e_0$}
\put(55,39){$\scriptstyle v_b$}
\put(32,26){$\scriptstyle e_1$}
\put(67,31){$\scriptstyle e_4$}
\put(77,22){$\scriptstyle v_a$}
\put(63,12){$\scriptstyle e_2$}
\put(83,12){$\scriptstyle e_3$}
\put(14,0){$\scriptstyle x_1$}
\put(59,0){$\scriptstyle x_2$}
\put(86,0){$\scriptstyle x_3$}
\end{picture}
\par\smallskip
{\renewcommand{\thefigure}{7.4.5}\refstepcounter{figure}
\normalfont\bfseries 图~\thefigure\label{fig:7-4-5}}
\end{minipage}\par\medskip

